\section{Kalimat Sempurna (Proposisi Tertutup)}

Seperti kepingan puzzle yang terserak berserakan pada Bab-bab pendahuluan di muka kelas, pemungkas segenap jerih-payah pembedahan Logika Predikat mengabdi pada satu mandat agung: **Mentransformasi kerancuan Variabel pada \textit{Kalimat Terbuka} hingga memadat mengental luruh menjadi aspal solid kepastian argumen Proposisi \textit{Tertutup}**. 

Ada 2 buah gerbang sakral untuk menyeberang menuju tahapan kalimat tertutup berkaliber valid, dan ketiadaan salah satunya menihilkan segalanya \cite{rosen2019}.

\subsection{Gerbang 1: Penanaman Eksekutor Konstanta Spesifik}

Jalan paling purba, instan, mendasar sekaligus radikal membunuh liarnya seekor Fungsi Proposisional bebas, tiada cara lain mematuhinya kecuali menangkapnya dan menghadiahi substitusi nama spesifik. Takkan ada lagi peluang interpretasi bercabang makna.

\begin{contoh}
Diberikan fungsi Proposisional Terbuka: \[ L(x, y) \equiv \text{''}x \text{ lari jauh lebih tangkas ketimbang } y\text{''} \]
(Kedua $x$ dan $y$ status luntang-lantung bebas. Fungsi ini belum bermakna sepeser pun secara biner B/S.)

Kita pilih mensubstitusikan target 1 saja:
\begin{multline*}
L(\text{Snoopy}, y) \equiv \text{``Anjing Snoopy berlari jauh lebih kencang}\\
\text{dari oknum } y\text{''}
\end{multline*}
(Fungsi malang ini ternyata masih setengah buta. Karena huruf $y$ masih berstatus liar, nilai kebenaran tetap urung tervalidasi).

Lampiaskan seluruh substitusi variabel tanpa sisa sebutirpun:
\begin{multline*}
L(\text{Snoopy}, \text{Kura-Kura Ninja}) \equiv \text{``Anjing Snoopy jauh lebih ngacir}\\
\text{ketimbang si Kura-kura Ninja''}
\end{multline*}
\textbf{Final Result:} Kalimat sukses tertutup! Tanpa ba-bi-bu argumen ini berhak mendapatkan seragam hakim diklasifikasikan \textbf{BENAR (T) atau SALAH (F)}, tergantung pada fakta alam satwa.
\end{contoh}

\subsection{Gerbang 2: Penjeratan Perangkap Kuantifikasi Makro}

Ada situasi miris bagi penjelajah di pedalaman matematika ketika mereka terpaksa menghadapi \textbf{lautan trilyunan kandidat tak bercita rasa semesta populasi} untuk disubstitusi. Sungguh melelahkan bila diwajibkan menjlentrehkan ratusan juta barisan $P(1), P(2), \ldots$ secara konstan berurut \cite{wiki_first_order_logic}. Solusinya? Logika Predikat mengerahkan teknologi agregat global massal: **KUANTIFIKASI**.

Dengan membubuhkan selubung pengikat jaring kuantor di wajah Variabel $x$ dan menyuntik syariat deklaratif di mukadimahnya, seluruh varians tak terhingga elemen acak langsung tersedot terkunci nilainya berbarengan menjadi komunal Proposisi tunggal masif.

\begin{contoh}
Tinjau kembali fungsi malang terbuka: 
\[ P(x) \equiv \text{''}x \text{ takkan hidup selamanya di galaksi ini\text{''}} \quad (\mathcal{U} = \text{biorganisme}) \]

Bermodalkan penambahan stempel magis Kuantor $\forall$ (Universal) serta $\exists$ (Eksistensial) pada bagian dada klausa:
\begin{enumerate}
    \item \textbf{Casting Kuantor Global}: $\forall x \, P(x)$ \\
    Meleburkan makna: \textit{"SELURUH biosistem yang pernah bernafas dari zaman sel primitif amuba hingga manusia modern, tanpa meleset walau sesosok tubuh pun, sejatinya takkan pernah berumur niscaya keabadian"}.
    Kalimat langsung terikat sakral sempurna $\implies$ Proposisi \textbf{Tertutup (BERNILAI BENAR T)}.

    \item \textbf{Casting Kuantor Mikro Tunggal}: $\exists x \, P(x)$ \\
    Bermakna hemat parsial moderat: \textit{"Setidaknya alam tak menampik ADA walau sebiji spesies saja kehidupan memilukan di alam semesta sana yang direnggut nyawanya sirna wafat ditelan era"}.
    Jaring ini juga sukses mengikat mutlak variabel anonim beribu spesies, menjadikan sebuah klausa argumen Proposisi \textbf{Tertutup super ringan (BENAR T)}.
\end{enumerate}
\end{contoh}

Jadi terbukti shahih dan niscaya valid, Logika Orde Pertama memberangus ketidakpastian Kalimat terbuka bukan sekedar dengan suap nama absolut entitas tunggal, namun didukung penuh komando agresi pembekuan sapu jagat massal bermodalkan \textbf{Pengunci Kuantifikasi $\forall$ dan $\exists$}. (Detail bahasan pengunci kuantifikasi dipelajari esok lusa pada palagan ranah Bab 7).
